\documentclass[11pt]{article}
\usepackage[margin=1in]{geometry}
\usepackage{booktabs}
\usepackage{amsmath}
\usepackage{enumitem}
\usepackage{longtable}

\title{Ablation-Guided Preset Selection for VieCut}
\author{}
\date{\today}

\begin{document}
\maketitle

\section{Reduction Set}
VieCut's min-cut reduction pipeline considered in this study consists of:
\begin{itemize}[leftmargin=*]
	\item \texttt{lp}: label-propagation contraction; merges locally similar vertices to shrink the graph quickly.
	
	\item \texttt{trivial}: local search refinement within clusters; identifies and moves ``mismatched" vertices out of small clusters
	
	\item \texttt{pr1}: Padberg--Rinaldi rule~1; contract $(u,v)$ if the edge dominates one endpoint, i.e.\ 
	$c(u,v) \ge \tfrac{1}{2} w(u)$ or $c(u,v) \ge \tfrac{1}{2} w(v)$.
	
	\item \texttt{pr2}: Padberg--Rinaldi rule~2; contract $(u,v)$ if the edge weight already exceeds a known lower bound on the global min-cut value, i.e.\ 
	$c(u,v) \ge \underline{\lambda}$.
	
	\item \texttt{pr3}: Padberg--Rinaldi rule~3; contract $(u,v)$ if it is not weaker than its connections through a common neighbor, i.e.\ 
	there exists $w$ with $c(u,v) \ge \min\{c(u,w),c(v,w)\}$.
	
	\item \texttt{pr4}: Padberg--Rinaldi rule~4; contract $(u,v)$ if separating the endpoints is no cheaper than isolating either vertex, i.e.\ 
	$c(u,v) \ge \min\{w(u)-c(u,v),\, w(v)-c(u,v)\}$.
\end{itemize}

\section{Class-to-Preset Ablation Summary}
Table~\ref{tab:class-preset-gains} summarizes, by graph class, the recommended preset from ablations, runtime gain over baseline, and change in pre-NOI structural reduction metrics.

\begin{table}[h]
\centering
\small
\begin{tabular}{l l r r r}
\toprule
Class & Preset & $\Delta t$ (s) & Speedup & $\Delta$ pre-NOI $n$ gain (\%) \\
\midrule
citation     & \texttt{no\_lp}      & 0.017 & 1.082 & 0.0 \\
co-purchase  & \texttt{no\_trivial} & 0.010 & 1.032 & 0.0 \\
delaunay     & \texttt{no\_trivial} & 0.011 & 1.174 & 0.0 \\
dyn-frames   & \texttt{no\_pr4}     & 1.260 & 1.087 & $-4.975\times10^5$ \\
er           & \texttt{no\_lp}      & 11.403 & 1.610 & $-3.146\times10^8$ \\
kronecker    & \texttt{no\_trivial} & 0.002 & 1.004 & 0.0 \\
matrix       & \texttt{no\_trivial} & 0.024 & 1.020 & 0.0 \\
mesh         & \texttt{no\_pr4}     & 0.020 & 1.009 & 0.0 \\
numerical    & \texttt{no\_trivial} & 0.509 & 1.145 & 0.0 \\
rgg          & \texttt{no\_trivial} & 0.007 & 1.019 & 0.0 \\
social       & \texttt{no\_pr4}     & 0.020 & 1.014 & 0.0 \\
street       & \texttt{no\_lp}      & 0.141 & 1.212 & 0.0 \\
web          & \texttt{no\_trivial} & 0.020 & 1.042 & 0.0 \\
\bottomrule
\end{tabular}
\caption{Preset recommendation by class.}
\label{tab:class-preset-gains}
\end{table}

\noindent\textbf{Cut correctness.} Cut value did not change in these experiments: all classes had zero cut mismatches vs baseline.

\section{High-Impact Classification Policy: \texttt{er}/\texttt{street} vs Rest}
From Table~\ref{tab:class-preset-gains}, the highest speedup impact is dominated by detecting \texttt{er} and \texttt{street} (both map to \texttt{no\_lp}); all remaining classes can safely use \texttt{no\_trivial} as the default.

Recommended policy:
\[
\{\texttt{er},\ \texttt{street},\ \texttt{other}\},\quad
\texttt{other} \rightarrow \texttt{no\_trivial}.
\]

Model (Histogram-based Gradient Boosting Classification Tree) quality for a compact 7-feature classifier (20 grouped holdout splits, grouped by graph):
\begin{itemize}[leftmargin=*]
    \item mean accuracy: $0.9690 \pm 0.0241$
    \item mean macro-F1: $0.9211 \pm 0.0640$
\end{itemize}

\subsection*{Top Choice Features (Minimal/Fast Set)}
\begin{enumerate}[leftmargin=*]
    \item \texttt{feature\_degree\_p50}: the median node degree (50th percentile). This is the ``typical'' connectivity level of a vertex and is robust to extreme hubs. If this value is high, the graph is globally more connected in a uniform sense; if low, the graph is generally sparse.
    \item \texttt{feature\_degree\_ratio\_max\_p90}: ratio of maximum degree to 90th-percentile degree. This measures how extreme the single largest hub is relative to already-high-degree vertices. Values near 1 indicate no exceptional hub; very large values indicate hub-dominated structure.
    \item \texttt{feature\_degree\_assortativity\_sampled}: sampled Pearson correlation of degrees at both endpoints of an edge. Positive values suggest high-degree nodes connect to high-degree nodes (assortative), while negative values indicate hub-to-low-degree linking (disassortative). This helps separate different random/network formation patterns.
    \item \texttt{feature\_kcore\_top\_fraction}: fraction of vertices that belong to the highest core index in a core decomposition. A large value means a substantial dense backbone; a small value means only a tiny elite core exists. This is useful for distinguishing uniform random-like connectivity from highly heterogeneous structure.
    \item \texttt{feature\_degree\_tail\_hill\_alpha}: Hill-estimator-style proxy of degree-tail shape using top positive degrees. Larger values correspond to a thinner tail (less extreme heavy-tail behavior), smaller values to heavier tails. This directly captures whether degree heterogeneity is mild or severe.
    \item \texttt{feature\_degree\_ratio\_p99\_p50}: ratio of 99th-percentile degree to median degree. This is a robust tail-vs-center spread indicator (less sensitive than max-based ratios). High values imply strong skew between typical and high-degree vertices.
    \item \texttt{feature\_degree\_p90}: 90th-percentile degree, i.e., the degree level exceeded by only 10\% of nodes. It summarizes the upper-but-not-extreme part of the distribution and stabilizes decisions that would otherwise depend too much on a single maximum-degree outlier.
\end{enumerate}

%\subsection*{Feature Computation Detail and Cost}
%\begin{table}[h]
%\centering
%\small
%\begin{tabular}{l l l}
%\toprule
%Feature & Computation (high level) & Cost tier \\
%\midrule
%\texttt{degree\_p50} & degree scan + quantile from sorted degree vector & cheap \\
%\texttt{degree\_ratio\_max\_p90} & max degree and p90 degree ratio & cheap \\
%\texttt{degree\_assortativity\_sampled} & sampled edge endpoint degree correlation & moderate sampled \\
%\texttt{kcore\_top\_fraction} & core decomposition + top-core count & medium \\
%\texttt{degree\_tail\_hill\_alpha} & Hill-style tail estimate on top positive degrees & cheap \\
%\texttt{degree\_ratio\_p99\_p50} & p99 to p50 degree ratio & cheap \\
%\texttt{degree\_p90} & degree quantile at 90th percentile & cheap \\
%\bottomrule
%\end{tabular}
%\caption{Selected 7-feature set with dominant computation used by the implementation.}
%\end{table}

%Most selected features are linear-pass degree statistics; the two heavier parts are sampled assortativity and one core decomposition pass.

%\subsection*{Measured Feature-Extraction Time}
%Feature extraction time (from \texttt{results/graph\_class\_guess\_more.csv}, using the current full extractor implementation) is:<
%\begin{itemize}[leftmargin=*]
%    \item median: $3.2067$ s
%    \item p90: $28.9123$ s
%    \item p95: $35.9276$ s
%\end{itemize}
%These values include computing the full feature superset currently emitted by \texttt{graph\_class\_guesser}, not only the selected 7-feature subset.

%\subsection*{Inference-Time Check}
%Measured median single-graph inference time for this 7-feature model is approximately
%\[
%4.43\ \text{ms}
%\]
%(in Python, model prediction only; excludes graph feature extraction).

%Since the deployed policy uses only the 7 features above (and omits BFS/clustering features), practical extraction overhead can be reduced further with a dedicated fast-path extractor.

%Even with this added inference overhead, the policy remains clearly beneficial relative to ablation gains:
%\begin{itemize}[leftmargin=*]
%    \item \texttt{er}: median runtime gain $11.403$ s
%    \item \texttt{street}: median runtime gain $0.141$ s
%\end{itemize}
%Thus, model inference time is negligible compared to the observed per-instance runtime savings.

\section{Dataset Used}
The dataset for this study contains 113 graph instances, including all DIMACS10 instances. The table below lists graph name, class label, and size $(n,m)$ for each instance.

\begingroup
\small
% Auto-generated from results/graph_class_guess_more.csv
\begin{longtable}{l l r r}
\toprule
Class & Graph name & $n$ & $m$ \\
\midrule
\endfirsthead
\toprule
Class & Graph name & $n$ & $m$ \\
\midrule
\endhead
\midrule
\multicolumn{4}{r}{\emph{Continued on next page}} \\
\midrule
\endfoot
\bottomrule
\endlastfoot
citation & citationCiteseer & 268495 & 1156647 \\
citation & coAuthorsCiteseer & 227320 & 814134 \\
citation & coAuthorsDBLP & 299067 & 977676 \\
citation & coPapersCiteseer & 434102 & 16036720 \\
citation & coPapersDBLP & 540486 & 15245729 \\
citation & dblp-2010 & 300647 & 807700 \\
co-purchase & amazon0312 & 400727 & 2349869 \\
co-purchase & amazon0505 & 410236 & 2439437 \\
co-purchase & amazon0601 & 403394 & 2443408 \\
delaunay & delaunay\_n10 & 1024 & 3056 \\
delaunay & delaunay\_n11 & 2048 & 6127 \\
delaunay & delaunay\_n12 & 4096 & 12264 \\
delaunay & delaunay\_n13 & 8192 & 24547 \\
delaunay & delaunay\_n14 & 16384 & 49122 \\
delaunay & delaunay\_n15 & 32768 & 98274 \\
delaunay & delaunay\_n16 & 65536 & 196575 \\
delaunay & delaunay\_n17 & 131072 & 393176 \\
delaunay & delaunay\_n18 & 262144 & 786396 \\
delaunay & delaunay\_n19 & 524288 & 1572823 \\
delaunay & delaunay\_n20 & 1048576 & 3145686 \\
delaunay & delaunay\_n21 & 2097152 & 6291408 \\
delaunay & delaunay\_n22 & 4194304 & 12582869 \\
delaunay & delaunay\_n23 & 8388608 & 25165784 \\
delaunay & delaunay\_n24 & 16777216 & 50331601 \\
dyn-frames & hugebubbles-00000 & 18318143 & 27470081 \\
dyn-frames & hugebubbles-00010 & 19458087 & 29179764 \\
dyn-frames & hugebubbles-00020 & 21198119 & 31790179 \\
dyn-frames & hugetrace-00000 & 4588484 & 6879133 \\
dyn-frames & hugetrace-00010 & 12057441 & 18082179 \\
dyn-frames & hugetrace-00020 & 16002413 & 23998813 \\
dyn-frames & hugetric-00000 & 5824554 & 8733523 \\
dyn-frames & hugetric-00010 & 6592765 & 9885854 \\
dyn-frames & hugetric-00020 & 7122792 & 10680777 \\
er & er-fact1.5-scale20 & 1048576 & 10904496 \\
er & er-fact1.5-scale21 & 2097152 & 22901945 \\
er & er-fact1.5-scale22 & 4194304 & 47983225 \\
er & er-fact1.5-scale23 & 8388608 & 100311279 \\
kronecker & kron\_g500-logn16 & 65536 & 3144264 \\
kronecker & kron\_g500-logn18 & 262144 & 12580729 \\
kronecker & kron\_g500-logn20 & 1048576 & 50328206 \\
kronecker & kron\_g500-logn21 & 2097152 & 100659103 \\
kronecker & kron\_g500-simple-logn16 & 65536 & 2456071 \\
kronecker & kron\_g500-simple-logn17 & 131072 & 5113985 \\
kronecker & kron\_g500-simple-logn18 & 262144 & 10582686 \\
kronecker & kron\_g500-simple-logn20 & 1048576 & 44619402 \\
kronecker & kron\_g500-simple-logn21 & 2097152 & 91040932 \\
matrix & G3\_circuit & 1585478 & 3037674 \\
matrix & af\_shell10 & 1508065 & 25582130 \\
matrix & af\_shell9 & 504855 & 8542010 \\
matrix & audikw1 & 943695 & 38354076 \\
matrix & cage15 & 5154859 & 47022346 \\
matrix & circuit5M & 5558311 & 26983926 \\
matrix & ecology1 & 1000000 & 1998000 \\
matrix & ecology2 & 999999 & 1997996 \\
matrix & hcircuit & 105676 & 203734 \\
matrix & kkt\_power & 2063494 & 6482320 \\
matrix & ldoor & 952203 & 22785136 \\
matrix & nlpkkt120 & 3542400 & 46651696 \\
matrix & nlpkkt160 & 8345600 & 110586256 \\
matrix & nlpkkt200 & 16240000 & 215992816 \\
matrix & nlpkkt240 & 27993600 & 373239376 \\
matrix & thermal2 & 1227087 & 3676134 \\
mesh & Bump\_2911 & 2852430 & 62409240 \\
mesh & Dubcova1.mtx-sorted & 16129 & 118440 \\
mesh & Flan\_1565 & 1564794 & 57920625 \\
mesh & HV15R & 2017169 & 162357569 \\
mesh & ML\_Laplace & 377002 & 13656485 \\
numerical & AS365 & 3799275 & 11368076 \\
numerical & M6 & 3501776 & 10501936 \\
numerical & NLR & 4163763 & 12487976 \\
numerical & adaptive & 6815744 & 13624320 \\
numerical & channel-500x100x100-b050 & 4802000 & 42681372 \\
numerical & nlr & 4163763 & 12487976 \\
numerical & packing-500x100x100-b050 & 2145852 & 17488243 \\
numerical & venturiLevel3 & 4026819 & 8054237 \\
rgg & rgg\_n\_2\_15\_s0 & 32768 & 160240 \\
rgg & rgg\_n\_2\_16\_s0 & 65536 & 342127 \\
rgg & rgg\_n\_2\_17\_s0 & 131072 & 728753 \\
rgg & rgg\_n\_2\_18\_s0 & 262144 & 1547283 \\
rgg & rgg\_n\_2\_19\_s0 & 524288 & 3269766 \\
rgg & rgg\_n\_2\_20\_s0 & 1048576 & 6891620 \\
rgg & rgg\_n\_2\_21\_s0 & 2097152 & 14487995 \\
rgg & rgg\_n\_2\_22\_s0 & 4194304 & 30359198 \\
rgg & rgg\_n\_2\_23\_s0 & 8388608 & 63501393 \\
rgg & rgg\_n\_2\_24\_s0 & 16777216 & 132557200 \\
social & com-amazon & 334863 & 925872 \\
social & com-dblp & 317080 & 1049866 \\
social & com-lj & 3997962 & 34681189 \\
social & com-orkut & 3072441 & 117185083 \\
social & com-youtube & 1134890 & 2987624 \\
social & ljournal-2008 & 5363186 & 49514271 \\
social & soc-LiveJournal1 & 4846609 & 42851237 \\
social & soc-flixster & 2523386 & 7918801 \\
social & soc-lastfm & 1191805 & 4519330 \\
social & soc-pokec & 1632803 & 22301964 \\
street & asia.osm & 11950757 & 12711603 \\
street & belgium.osm & 1441295 & 1549970 \\
street & ca-hollywood-2009 & 1069126 & 56306653 \\
street & europe.osm & 50912018 & 54054660 \\
street & germany.osm & 11548845 & 12369181 \\
street & great-britain.osm & 7733822 & 8156517 \\
street & italy.osm & 6686493 & 7013978 \\
street & luxembourg.osm & 114599 & 119666 \\
street & netherlands.osm & 2216688 & 2441238 \\
street & roadNet-CA & 1965206 & 2766607 \\
street & roadNet-PA & 1088092 & 1541898 \\
street & roadNet-TX & 1379917 & 1921660 \\
web & eu-2005 & 862664 & 16138468 \\
web & in-2004 & 1382908 & 13591473 \\
web & web-Google-sorted & 356648 & 2093324 \\
web & web-NotreDame & 325729 & 1090108 \\
web & wiki-Talk & 232314 & 1458806 \\
web & wiki-topcats & 1791489 & 25444207 \\
\end{longtable}

\endgroup

\end{document}
